Organic solar cells have been studied since x when y started researching them.
Progress in OSC develop consistently lagged behind silicon cells, but the envention of Y6:PM6 system suddenly launched the effieciency of OSCs in \textcolor{red}{year}.


\subsection{Advantages of OSCs}
A promising feature of OSCs is the ability for them to be manufactored efficiently on a rall to roll process.

\subsection{Challenges with OSCs}
Historically, OSCs have been plagued with low PCE.
The 10\% PCE mark was broken in \textcolor{red}{year} and 20\% PCE was recently broken in \textcolor{red}{year}.
A leading challenge in the development of practical OSCs is the lifespan of the cells.
Lifespan is measured at T\textsubscript{80}, the time it takes for degredation of the eficiency has lead the product to only be 80\% the eficiency o the new product\autocite{duEfficientPolymerSolar2019}.
There are multiple causes for the degradation of OSC mateials, including moisure, heat, and irradiateion\autocite{tetreaultAssessingIndividualMaterial2024}.

\subsection{Stille coupling reaction}
The stille coupling reaction is ubiquitous in the synthesis of materials for OSCs.
An alkylated tin moiety can be coupled to an aromatic halide to form a conjugated pi system.
This capability is powerfull, as many NFAs and plymeric donor materials are highly conjugated.
Commonly, polymeric or oligomeric conjugated thiophenes are used to build the acceptors and donors of the BHJ.

\subsection{Previous work}
There have been attempts to computaionally screen materials for use in OSCs.
